\chapter{الاستبدال المحب للنواة}\label{ch:b1:nucleophilic-substitution}

خذ قارورة من 2-بروموبوتان-\cip{S}، وهو \omterm{def:b1:stereochemistry-in-depth:enantiomers}{متماكب مرآوي} وحيد يدوّر
الضوء المستقطب. فإذا سُخّن مع هيدروكسيد الصوديوم في \omterm{def:b1:intermolecular-forces-solvents:solvent-classes}{مذيب قطبي} لابروتوني،
أعطى بوتان-2-أول هو أيضًا \omterm{def:b1:stereochemistry-in-depth:enantiomers}{متماكب مرآوي} وحيد --- لكنه
الآخر، \cip{R}: لقد قلب التفاعل الجزيء كما تقلب الريح
مظلة. أما إذا تُرك البروميد نفسه في ماء دافئ، فإنه يعطي
بوتان-2-أول راسيميًا جزئيًا. التحول الإجمالي نفسه، مجموعة OH
تحل محل ذرة Br، يسلك مسارين مختلفين. يثبت هذا الفصل
الآليتين انطلاقًا من قانوني سرعتهما، ويشرح
كيميائهما الفراغية، ويعطي القواعد التي تحسم بينهما --- وهو أول
استعمال كامل لأدوات فصلي الحركية والتأثيرات
الإلكترونية.

\begin{recall}
\omterm{def:b1:rate-laws:order}{قوانين السرعة}، \omterm{def:b1:elementary-steps:elementary-step}{والخطوات الأولية}، \omterm{def:b1:elementary-steps:rds}{والخطوة المحددة للسرعة}
\omterm{def:b1:elementary-steps:qssa}{وتقريب الحالة المستقرة} موجودة في \cref{ch:b1:rate-laws,ch:b1:elementary-steps}؛
والترتيبات الفراغية، وتمثيلات كرام \omterm{def:b1:stereochemistry-in-depth:optical-activity}{والنشاط الضوئي} في
\cref{ch:b1:stereochemistry-in-depth}؛ \omterm{def:b1:electronic-effects:nucleophile}{والمحبات للنواة}، \omterm{def:b1:electronic-effects:nucleophile}{والمحبات للإلكترونات}،
\omterm{def:b1:electronic-effects:nucleophile}{والمجموعات المغادرة} \omterm{def:b1:electronic-effects:intermediates}{والكاتيونات الكربونية} في \cref{ch:b1:electronic-effects}؛
\omterm{def:b1:intermolecular-forces-solvents:solvent-classes}{والمذيبات البروتونية} واللابروتونية في \cref{ch:b1:intermolecular-forces-solvents}.
\end{recall}

\section{الاستبدال وأطرافه}

\begin{definition}[الاستبدال المحب للنواة، الركيزة]\label{def:b1:nucleophilic-substitution:substitution}
\emph{الاستبدال المحب للنواة}\index{استبدال محب للنواة} هو
إحلال \omterm{def:b1:electronic-effects:nucleophile}{محب للنواة} Nu، على كربون مشبع، محل \omterm{def:b1:electronic-effects:nucleophile}{مجموعة مغادرة} Y:
$\mathrm{Nu^-} + \mathrm{R{-}Y} \to \mathrm{R{-}Nu} +
\mathrm{Y^-}$. والمركب R--Y هو \emph{الركيزة}\index{ركيزة}؛
والكربون الحامل للمجموعة Y \omterm{def:b1:electronic-effects:nucleophile}{محب للإلكترونات} لأن الرابطة C--Y مستقطبة
نحو Y.
\end{definition}

الهالوجينو ألكانات هي الركائز النموذجية، ومجموعاتها المغادرة \ce{Cl-} أو \ce{Br-} أو
\ce{I-}. فالهيدروكسيد يعطي كحولات؛ والألكوكسيدات تعطي
إيثرات؛ والسيانيد يعطي نتريلات (بكربون إضافي)؛ والأمونيا والأمينات تعطي
أمينات؛ واليوديد يتبادل مع الهاليدات الأخرى. ويمكن للماء والكحولات، حين تُستعمل
مذيبات، أن تؤدي أيضًا دور \omterm{def:b1:electronic-effects:nucleophile}{المحبات للنواة}: فيكون الاستبدال عندئذ
تحللًا بالمذيب.

\section{الآلية SN2}

حين يتفاعل البروموميثان مع الهيدروكسيد في الماء، فإن مضاعفة أي من
التركيزين تضاعف السرعة: $v = k[\ce{CH3Br}][\ce{OH-}]$، أي تفاعل من
الرتبة الأولى بالنسبة إلى كل متفاعل.

\begin{definition}[الآلية SN2، انقلاب والدن]\label{def:b1:nucleophilic-substitution:sn2}
\emph{الآلية SN2}\index{آلية SN2} (\omterm{def:b1:nucleophilic-substitution:substitution}{استبدال
محب للنواة} ثنائي الجزيئية) \omterm{def:b1:elementary-steps:elementary-step}{خطوة أولية} واحدة يهاجم فيها
\omterm{def:b1:electronic-effects:nucleophile}{المحب للنواة} الكربون من الجهة المقابلة \omterm{def:b1:electronic-effects:nucleophile}{للمجموعة المغادرة}
بينما تغادر \omterm{def:b1:electronic-effects:nucleophile}{المجموعة المغادرة}. وتنقلب متبادلات الكربون الثلاثة
الأخرى، كمظلة قلبتها الريح: وهذا \emph{انقلاب
والدن}\index{انقلاب والدن}.
\end{definition}

\begin{omfigure}
\setchemfig{atom sep=2.1em}
\schemestart
\chemfig{H@{nu}O^{-}}
\arrow{0}[,0.35]
\chemfig{@{c}C(-[:150]H_3C)(<[:210]H)(<:[:250]C_2H_5)-[@{cb}]@{br}Br}
\arrow{->}
\chemfig{[:0]HO-C(-[:30]CH_3)(<[:-30]H)(<:[:-70]C_2H_5)}
\+
\chemfig{Br^{-}}
\schemestop
\chemmove{\draw[->, thick, shorten <=2pt, shorten >=3pt](nu) .. controls +(15:5mm) and +(175:5mm) .. (c);
          \draw[->, thick, shorten <=2pt](cb) .. controls +(90:6mm) and +(90:6mm) .. (br);}

\medskip
\begin{tikzpicture}[font=\small]
  \node[draw, dashed, rounded corners, inner sep=6pt] at (0,0)
    {$\left[\vcenter{\hbox{\ce{HO}$\cdots$\chemfig{C(-[2]CH_3)(-[:-30]H)(-[:210]C_2H_5)}$\cdots$\ce{Br}}}\right]^{-}$};
  \node[below] at (0,-1.1) {\foreignlanguage{arabic}{الحالة الانتقالية: خمس مجموعات حول \babelsublr{C}، ثلاث منها في مستوٍ}};
\end{tikzpicture}

{\small التفاعل SN2 للهيدروكسيد مع 2-بروموبوتان-\cip{S}. يهاجم
\omterm{def:b1:electronic-effects:nucleophile}{المحب للنواة} الوجه المقابل للبروم (\omterm{def:b1:electronic-effects:curly-arrow}{الأسهم المنحنية})؛ وفي
\omterm{def:b1:elementary-steps:energy-profile}{الحالة الانتقالية} تتكوّن الرابطة O--C بينما تنكسر الرابطة C--Br؛
والناتج بوتان-2-أول-\cip{R}: فقد انقلبت المجموعات الثلاث الأخرى إلى
الجهة الأخرى.}
\end{omfigure}

\begin{proposition}[قانون سرعة SN2]\label{prop:b1:nucleophilic-substitution:sn2-law}
من أجل تفاعل SN2، $v = k[\mathrm{RY}][\mathrm{Nu}]$.
\end{proposition}

\begin{proof}
الآلية \omterm{def:b1:elementary-steps:elementary-step}{خطوة أولية} واحدة ثنائية الجزيئية؛ وبقانون فعل الكتلة
\omterm{def:b1:elementary-steps:elementary-step}{للخطوات الأولية} (\cref{ch:b1:elementary-steps}) تكون سرعتها
متناسبة مع جداء تركيزي الطرفين.
\end{proof}

\begin{proposition}[الكيمياء الفراغية للآلية SN2]\label{prop:b1:nucleophilic-substitution:sn2-inversion}
يعطي التفاعل SN2 على كربون فراغي المركز \omterm{def:b1:stereochemistry-in-depth:isomers}{متماكبًا فراغيًا} وحيدًا، يكون فيه
ترتيب المجموعات الثلاث غير المتغيرة مقلوبًا: فالتفاعل من هذا
النوع نوعي فراغيًا.
\end{proposition}

\begin{proof}
لا يستطيع \omterm{def:b1:electronic-effects:nucleophile}{المحب للنواة} أن يبلغ إلا الجهة الخالية من الكربون، المقابلة
\omterm{def:b1:electronic-effects:nucleophile}{للمجموعة المغادرة}؛ فتعبر المجموعات الثلاث الأخرى مستويًا في
\omterm{def:b1:elementary-steps:energy-profile}{الحالة الانتقالية} وتنتهي في الجهة الأخرى. وتتجه الرابطة الجديدة حيث
لم تكن القديمة: فكل جزيء من \omterm{def:b1:nucleophilic-substitution:substitution}{الركيزة} يعطي الترتيب المرآوي
لمجموعاته. ويتغيّر الواصف عادةً من \cip{S}
إلى \cip{R} أو العكس، لكن ليس دائمًا، لأنه يتعلق برتبتي
\omterm{def:b1:electronic-effects:nucleophile}{المحب للنواة} \omterm{def:b1:electronic-effects:nucleophile}{والمجموعة المغادرة}.
\end{proof}

\begin{definition}[التفاعل النوعي فراغيًا]\label{def:b1:nucleophilic-substitution:stereospecific}
يكون التفاعل \emph{تفاعلًا نوعيًا فراغيًا}\index{تفاعل نوعي فراغيًا}
حين تعطي ركائزُ متماكبة فراغيًا نواتجَ مختلفة
فراغيًا، فتعطي كل \omterm{def:b1:nucleophilic-substitution:substitution}{ركيزة} ناتجها الخاص.
\end{definition}

\section{الآلية SN1}

يتفاعل 2-برومو-2-ميثيل البروبان مع الماء بسرعة لا تتعلق
بتركيز \omterm{def:b1:electronic-effects:nucleophile}{المحب للنواة}: $v = k[(\ce{CH3})_3\ce{CBr}]$. فلا بد أن
خطوة لا يشارك فيها \omterm{def:b1:electronic-effects:nucleophile}{المحب للنواة} هي التي تحدد السرعة.

\begin{definition}[الآلية SN1، التحول الراسيمي، التحلل بالمذيب]\label{def:b1:nucleophilic-substitution:sn1}
\emph{الآلية SN1}\index{آلية SN1} (\omterm{def:b1:nucleophilic-substitution:substitution}{استبدال
محب للنواة} أحادي الجزيئية) لها خطوتان: \omterm{def:b1:electronic-effects:cleavage}{الانشطار غير المتجانس} البطيء للرابطة
C--Y، الذي يعطي \omterm{def:b1:electronic-effects:intermediates}{كاتيونًا كربونيًا}، ثم الإضافة السريعة
\omterm{def:b1:electronic-effects:nucleophile}{للمحب للنواة} إلى \omterm{def:b1:electronic-effects:intermediates}{الكاتيون الكربوني}. والتفاعل الذي يحوّل \omterm{def:b1:stereochemistry-in-depth:enantiomers}{متماكبًا مرآويًا} وحيدًا
إلى خليط من الاثنين \emph{تحول راسيمي}\index{تحول راسيمي}؛
والاستبدال الذي يكون فيه المذيب هو \omterm{def:b1:electronic-effects:nucleophile}{المحب للنواة}
\emph{تحلل بالمذيب}\index{تحلل بالمذيب}.
\end{definition}

\begin{omfigure}
\setchemfig{atom sep=2.1em}
\schemestart
\chemfig{C(-[2]CH_3)(-[4]H_3C)(-[6]CH_3)-[@{cb2}]@{br2}Br}
\arrow{->[بطيئة]}
\chemfig{C^{+}(-[2]CH_3)(-[:210]H_3C)(-[:-30]CH_3)}
\+
\chemfig{Br^{-}}
\schemestop

\smallskip
\schemestart
\chemfig{C^{+}(-[2]CH_3)(-[:210]H_3C)(-[:-30]CH_3)}
\arrow{->[سريعة][\ce{H2O}]}
\chemfig{C(-[2]CH_3)(-[4]H_3C)(-[6]CH_3)-OH_2^{+}}
\arrow{->[\ce{-H+}]}
\chemfig{C(-[2]CH_3)(-[4]H_3C)(-[6]CH_3)-OH}
\schemestop
\chemmove{\draw[->, thick, shorten <=2pt](cb2) .. controls +(90:6mm) and +(90:6mm) .. (br2);}

\medskip
\begin{tikzpicture}[font=\small]
  \fill[omProp!15, draw=omProp] (0,0.2) ellipse (0.2 and 0.85);
  \fill[omProp!15, draw=omProp] (0,-0.2) ellipse (0.2 and 0.85);
  \draw[thick] (0,0) -- (-1.1,-0.2) node[left] {\ce{CH3}};
  \draw[thick] (0,0) -- (0.95,-0.45) node[right] {\ce{C2H5}};
  \draw[thick] (0,0) -- (0.4,0.5) node[above right] {H};
  \node[fill=white, inner sep=1pt] at (0,0) {\ce{C+}};
  \draw[->, thick, omDef] (-0.9,1.6) node[left] {\ce{H2O}} -- (-0.15,0.95);
  \draw[->, thick, omDef] (0.9,-1.6) node[right] {\ce{H2O}} -- (0.15,-0.95);
  \node[align=center] at (4.6,0) {\foreignlanguage{arabic}{هجوم على أيٍّ من الوجهين:}\\ \foreignlanguage{arabic}{الترتيبان كلاهما،}\\ وخليط راسيمي إن لم يختلف\\ شيء آخر بين الوجهين};
\end{tikzpicture}

{\small \omterm{def:b1:nucleophilic-substitution:sn1}{التحلل بالمذيب} SN1 للمركب 2-برومو-2-ميثيل البروبان في الماء: تأين بطيء
إلى \omterm{def:b1:electronic-effects:intermediates}{الكاتيون الكربوني} الثالثي، ثم أسر الماء له وفقد بروتون.
في الأسفل: \omterm{def:b1:electronic-effects:intermediates}{الكاتيون الكربوني} الثانوي \omterm{def:b1:stereochemistry-in-depth:stereocentre}{الكيرالي} (من 2-بروموبوتان) مستوٍ،
ويضاف الماء على أيٍّ من وجهيه باحتمال متساوٍ.}
\end{omfigure}

\begin{proposition}[قانون سرعة SN1]\label{prop:b1:nucleophilic-substitution:sn1-law}
من أجل \omterm{def:b1:nucleophilic-substitution:sn1}{الآلية SN1}، إذا أُسر \omterm{def:b1:electronic-effects:intermediates}{الكاتيون الكربوني} أسرع بكثير مما
يعود إلى \omterm{def:b1:nucleophilic-substitution:substitution}{الركيزة}، فإن $v = k_1[\mathrm{RY}]$، مستقلة عن
\omterm{def:b1:electronic-effects:nucleophile}{المحب للنواة}.
\end{proposition}

\begin{proof}
الخطوات: $\mathrm{RY} \rightleftharpoons \mathrm{R^+} + \mathrm{Y^-}$ ($k_1$،
$k_{-1}$)، $\mathrm{R^+} + \mathrm{Nu} \to \mathrm{RNu}$ ($k_2$).
\omterm{def:b1:electronic-effects:intermediates}{والكاتيون الكربوني} \omterm{def:b1:electronic-effects:intermediates}{وسيط تفاعلي}: فبتقريب الحالة
المستقرة (\cref{ch:b1:elementary-steps})، $k_1[\mathrm{RY}] =
k_{-1}[\mathrm{R^+}][\mathrm{Y^-}] + k_2[\mathrm{R^+}][\mathrm{Nu}]$، ومن ثم
$v = k_2[\mathrm{R^+}][\mathrm{Nu}] = \dfrac{k_1k_2[\mathrm{RY}][\mathrm{Nu}]}
{k_{-1}[\mathrm{Y^-}] + k_2[\mathrm{Nu}]}$. وحين يكون $k_2[\mathrm{Nu}] \gg
k_{-1}[\mathrm{Y^-}]$ يؤول هذا إلى $k_1[\mathrm{RY}]$: فالخطوة الأولى هي
\omterm{def:b1:elementary-steps:rds}{الخطوة المحددة للسرعة}.
\end{proof}

\begin{proposition}[الكيمياء الفراغية للآلية SN1]\label{prop:b1:nucleophilic-substitution:sn1-stereo}
يعطي التفاعل SN1 على كربون فراغي المركز الترتيبين كليهما: فالتفاعل
ليس نوعيًا فراغيًا، ويؤدي إلى \omterm{def:b1:nucleophilic-substitution:sn1}{تحول راسيمي} حين لا يميّز شيء
بين وجهي \omterm{def:b1:electronic-effects:intermediates}{الكاتيون الكربوني}.
\end{proposition}

\begin{proof}
\omterm{def:b1:electronic-effects:intermediates}{الكاتيون الكربوني} مستوٍ (\cref{ch:b1:electronic-effects}): فوجهاه
صورتان مرآويتان. وإذا ابتعدت \omterm{def:b1:electronic-effects:nucleophile}{المجموعة المغادرة} كثيرًا قبل
الأسر، أضيف \omterm{def:b1:electronic-effects:nucleophile}{المحب للنواة} إلى أيٍّ من الوجهين بالسرعة نفسها، فيعطي
\omterm{def:b1:stereochemistry-in-depth:enantiomers}{المتماكبين المرآويين} بكميتين متساويتين. وفي الواقع يبقى الأيون المغادر
حاجبًا لأحد الوجهين برهةً، وكثيرًا ما يُرصد فائض طفيف من الانقلاب.
\end{proof}

\begin{omfigure}
\begin{tikzpicture}
\begin{axis}[
  width=0.8\linewidth, height=5.0cm, xmin=0, xmax=10, ymin=-30, ymax=110,
  axis x line=bottom, axis y line=left, xlabel={إحداثي التفاعل},
  ylabel={الطاقة}, xtick=\empty, ytick=\empty,
  label style={font=\small}, legend style={font=\scriptsize, draw=none, at={(0.98,0.98)}, anchor=north east}]
\addplot[omDef, very thick] table[x=x, y=E] {figdata/out/bachelor-1-energy-profiles-sn2.dat};
\addplot[omProp, very thick, dashed] table[x=x, y=E] {figdata/out/bachelor-1-energy-profiles-sn1.dat};
\legend{{\foreignlanguage{arabic}{\babelsublr{SN2}: خطوة واحدة}}, {\foreignlanguage{arabic}{\babelsublr{SN1}: خطوتان}}}
\node[font=\scriptsize, anchor=north] at (axis cs:5,68) {كاتيون كربوني};
\node[font=\scriptsize, anchor=north] at (axis cs:0.6,-3) {RY};
\node[font=\scriptsize, anchor=south] at (axis cs:9.6,-18) {RNu};
\end{axis}
\end{tikzpicture}

{\small منحنيات طاقة تخطيطية. SN2: \omterm{def:b1:elementary-steps:energy-profile}{حالة انتقالية} واحدة، خمس
مجموعات حول الكربون. SN1: حاجز أول أعلى (التأين، \omterm{def:b1:elementary-steps:rds}{الخطوة
المحددة للسرعة})، ثم \omterm{def:b1:electronic-effects:intermediates}{الكاتيون الكربوني} في بئر ضحلة، ثم حاجز
صغير لأسره.}
\end{omfigure}

\section{ما الذي يحسم الآلية}

\begin{proposition}[الركيزة]\label{prop:b1:nucleophilic-substitution:substrate}
تتفاعل الركائز الميثيلية والأولية \omterm{def:b1:nucleophilic-substitution:sn2}{بالآلية SN2}؛ والركائز الثالثية \omterm{def:b1:nucleophilic-substitution:sn1}{بالآلية SN1}؛
والركائز الثانوية بأيٍّ منهما، تبعًا للعوامل الأخرى. وتتفاعل الركائز الأليلية
والبنزيلية بسرعة بالآليتين كلتيهما.
\end{proposition}

\begin{proof}
تحتاج SN2 إلى متسع خلف الكربون: فكل مجموعة ألكيل تحل محل هيدروجين
تزحم \omterm{def:b1:elementary-steps:energy-profile}{الحالة الانتقالية}، التي تحمل خمس مجموعات، وتبطئ
التفاعل؛ وثلاث مجموعات تسدّه. وتحتاج SN1 إلى \omterm{def:b1:electronic-effects:intermediates}{كاتيون كربوني} ثابت:
فترتيب الثبات (\cref{prop:b1:electronic-effects:carbocation-order})
يجعل التأين سريعًا للثالثي، وممكنًا للثانوي، وبطيئًا جدًا
\omterm{def:b1:electronic-effects:intermediates}{للكاتيونات الكربونية} الأولية والميثيلية. \omterm{def:b1:electronic-effects:intermediates}{والكاتيونات الكربونية} الأليلية والبنزيلية
تثبت بالرنين، والنظام $\pi$ نفسه يثبّت \omterm{def:b1:elementary-steps:energy-profile}{الحالة الانتقالية}
\omterm{def:b1:nucleophilic-substitution:sn2}{للآلية SN2}.
\end{proof}

\begin{proposition}[المجموعة المغادرة]\label{prop:b1:nucleophilic-substitution:leaving-group}
تكون الآليتان كلتاهما أسرع كلما كانت \omterm{def:b1:electronic-effects:nucleophile}{المجموعة المغادرة} أفضل، أي كلما كانت
القاعدة التي تصير إليها أضعف: $\ce{I-} > \ce{Br-} > \ce{Cl-} \gg \ce{F-}$؛
أما \ce{OH-} و\ce{RO-} فلا تغادران.
\end{proposition}

\begin{proof}
في الآليتين كلتيهما تنكسر الرابطة C--Y في \omterm{def:b1:elementary-steps:rds}{الخطوة المحددة للسرعة} أو قبلها،
وتحمل \omterm{def:b1:elementary-steps:energy-profile}{الحالة الانتقالية} جزءًا من الشحنة السالبة على Y.
والمجموعة التي تحمل الشحنة السالبة جيدًا --- القاعدة المرافقة \omterm{def:b1:predominant-reaction:strong-acid}{لحمض
قوي} --- تخفض تلك \omterm{def:b1:elementary-steps:energy-profile}{الحالة الانتقالية}. و\ce{HI} و\ce{HBr}
و\ce{HCl} \omterm{def:b1:predominant-reaction:strong-acid}{أحماض قوية}، و\ce{HF} \omterm{def:b1:predominant-reaction:strong-acid}{حمض ضعيف} ($\mathrm{p}K_a$ يساوي 3.2)، والماء
والكحولات \omterm{def:b1:predominant-reaction:strong-acid}{أحماض ضعيفة} جدًا: فالأيونان \ce{OH-} و\ce{RO-} قاعدتان قويتان
ومجموعتان مغادرتان رديئتان.
\end{proof}

\begin{proposition}[المذيب]\label{prop:b1:nucleophilic-substitution:solvent}
\omterm{def:b1:intermolecular-forces-solvents:solvent-classes}{المذيبات القطبية} البروتونية (الماء، الكحولات) تحابي SN1؛ \omterm{def:b1:intermolecular-forces-solvents:solvent-classes}{والمذيبات القطبية} اللابروتونية
(البروبانون، وثنائي ميثيل السلفوكسيد، وثنائي ميثيل الفورماميد) تحابي SN2 مع \omterm{def:b1:electronic-effects:nucleophile}{المحبات للنواة}
الأنيونية.
\end{proposition}

\begin{proof}
تكوّن SN1 أيونين من \omterm{def:b1:nucleophilic-substitution:substitution}{ركيزة} متعادلة: \omterm{def:b1:intermolecular-forces-solvents:solvent-classes}{والمذيب القطبي} البروتوني
يثبّت كليهما، الكاتيون بأزواجه الحرة والأنيون \omterm{def:b1:intermolecular-forces-solvents:hbond}{بالروابط
الهيدروجينية}، فيخفض حاجز التأين. وفي SN2 تنتشر شحنة
\omterm{def:b1:electronic-effects:nucleophile}{المحب للنواة} الأنيوني في \omterm{def:b1:elementary-steps:energy-profile}{الحالة الانتقالية}؛ \omterm{def:b1:intermolecular-forces-solvents:solvent-classes}{والمذيب البروتوني}،
الذي يرتبط بالأنيون الصغير بقوة \omterm{def:b1:intermolecular-forces-solvents:hbond}{بالروابط الهيدروجينية}، يبطئه،
بينما يترك \omterm{def:b1:intermolecular-forces-solvents:solvent-classes}{المذيب اللابروتوني} الأنيون حرًا ونشطًا
(\cref{ch:b1:intermolecular-forces-solvents}).
\end{proof}

\begin{method}[اختيار الآلية]\label{met:b1:nucleophilic-substitution:choose-mechanism}
\begin{enumerate}
  \item \omterm{def:b1:nucleophilic-substitution:substitution}{الركيزة}: \babelsublr{\foreignlanguage{arabic}{ميثيلية أو أولية} $\to$ SN2}؛ \babelsublr{\foreignlanguage{arabic}{ثالثية} $\to$ SN1}؛
        \babelsublr{\foreignlanguage{arabic}{ثانوية} $\to$ \foreignlanguage{arabic}{تابِع}}.
  \item \omterm{def:b1:electronic-effects:nucleophile}{المحب للنواة}: \babelsublr{\foreignlanguage{arabic}{قوي وأنيوني ($\ce{OH-}$، $\ce{RO-}$، $\ce{CN-}$،
        $\ce{I-}$)} $\to$ SN2}؛ \babelsublr{\foreignlanguage{arabic}{ضعيف ومتعادل (الماء، الكحول)} $\to$ SN1}.
  \item المذيب: \babelsublr{\foreignlanguage{arabic}{لابروتوني} $\to$ SN2}؛ \babelsublr{\foreignlanguage{arabic}{بروتوني} $\to$ SN1}.
  \item توقّع الكيمياء الفراغية (انقلاب أو \omterm{def:b1:nucleophilic-substitution:sn1}{تحول راسيمي}) وتحقق
        من أن \omterm{def:b1:predominant-reaction:strong-acid}{قاعدة قوية} لا تسبب بالأحرى حذفًا
        (\cref{ch:b1:elimination}).
\end{enumerate}
\end{method}

\begin{omfigure}
\begin{tabular}{lll}
\hline
 & SN2 & SN1 \\
\hline
الخطوات & واحدة & اثنتان، عبر \omterm{def:b1:electronic-effects:intermediates}{كاتيون كربوني} \\
\omterm{def:b1:rate-laws:order}{قانون السرعة} & $k[\mathrm{RY}][\mathrm{Nu}]$ & $k[\mathrm{RY}]$ \\
\omterm{def:b1:nucleophilic-substitution:substitution}{الركيزة} & \babelsublr{\foreignlanguage{arabic}{ميثيلية} $>$ \foreignlanguage{arabic}{أولية} $>$ \foreignlanguage{arabic}{ثانوية}} & \babelsublr{\foreignlanguage{arabic}{ثالثية} $>$ \foreignlanguage{arabic}{ثانوية}} \\
\omterm{def:b1:electronic-effects:nucleophile}{المحب للنواة} & قوي، أنيوني & ضعيف، متعادل (غالبًا المذيب) \\
المذيب & قطبي لابروتوني & قطبي بروتوني \\
الكيمياء الفراغية & انقلاب (نوعي فراغيًا) & \omterm{def:b1:nucleophilic-substitution:sn1}{تحول راسيمي} (غالبًا) \\
\hline
\end{tabular}

\smallskip
{\small الآليتان جنبًا إلى جنب.}
\end{omfigure}

\section{تمارين}

\begin{exercise}[$\star$]\label{exo:b1:nucleophilic-substitution:1}
اكتب نواتج: البروموإيثان مع هيدروكسيد الصوديوم؛ واليودوميثان مع
إيثوكسيد الصوديوم؛ و1-كلوروبروبان مع سيانيد البوتاسيوم؛ والبروموميثان مع
الأمونيا (الخطوة الأولى). وعيّن \omterm{def:b1:nucleophilic-substitution:substitution}{الركيزة} \omterm{def:b1:electronic-effects:nucleophile}{والمحب للنواة} \omterm{def:b1:electronic-effects:nucleophile}{والمجموعة المغادرة}.
\end{exercise}

\begin{exercise}[$\star$]\label{exo:b1:nucleophilic-substitution:2}
من أجل تفاعل $\mathrm{RY} + \mathrm{Nu}$ له $v = k[\mathrm{RY}][\mathrm{Nu}]$،
ماذا يحدث \omterm{def:b1:rate-laws:initial-rate}{للسرعة الابتدائية} إذا ضوعف $[\mathrm{Nu}]$ ثلاث مرات؟ وإذا نُصّف
التركيزان كلاهما؟ والأسئلة نفسها إذا كان $v = k[\mathrm{RY}]$.
\end{exercise}

\begin{exercise}[$\star$]\label{exo:b1:nucleophilic-substitution:3}
توقّع الآلية (SN1 أو SN2): 1-بروموبوتان مع يوديد الصوديوم في
البروبانون؛ 2-برومو-2-ميثيل البروبان في الماء؛ البروموميثان مع الهيدروكسيد؛
2-بروموبوتان في الإيثانول دون إضافة قاعدة.
\end{exercise}

\begin{exercise}[$\star$]\label{exo:b1:nucleophilic-substitution:4}
ارسم ناتج التفاعل SN2 للسيانيد مع 2-بروموبوتان-\cip{R}
\omterm{def:b1:stereochemistry-in-depth:representations}{بتمثيل كرام} وأعطِ واصفه.
\end{exercise}

\begin{exercise}[$\star\star$]\label{exo:b1:nucleophilic-substitution:5}
اكتب \omterm{def:b1:nucleophilic-substitution:sn1}{الآلية SN1} الكاملة للتحلل المائي للمركب 2-برومو-2-ميثيل البروبان،
\omterm{def:b1:electronic-effects:curly-arrow}{بالأسهم المنحنية}، بما في ذلك فقد البروتون. لماذا كان \omterm{def:b1:rate-laws:order}{قانون السرعة}
من الرتبة الأولى مع أن ثلاثة أنواع كيميائية تشارك فيه؟
\end{exercise}

\begin{exercise}[$\star\star$]\label{exo:b1:nucleophilic-substitution:6}
فسّر لماذا يتفاعل 1-برومو-2,2-ثنائي ميثيل البروبان، وهو بروميد أولي، ببطء
شديد \omterm{def:b1:nucleophilic-substitution:sn2}{بالآلية SN2} ولا يتفاعل \omterm{def:b1:nucleophilic-substitution:sn1}{بالآلية SN1}.
\end{exercise}

\begin{exercise}[$\star\star$]\label{exo:b1:nucleophilic-substitution:7}
يفقد 2-يودو أوكتان-\cip{S} نشاطه الضوئي حين يُحفظ في البروبانون
مع يوديد الصوديوم، مع أنه لا يتكوّن أي مركب جديد. فسّر ذلك، وبيّن لماذا
كانت سرعة فقد \omterm{def:b1:stereochemistry-in-depth:optical-activity}{النشاط الضوئي} ضعف سرعة الاستبدال.
\end{exercise}

\begin{exercise}[$\star\star$]\label{exo:b1:nucleophilic-substitution:8}
رتّب بحسب سرعة SN2 مع اليوديد: 2-برومو-2-ميثيل البروبان، والبروموميثان،
و2-بروموبروبان، و1-بروموبروبان. ثم رتّب المركبات نفسها بحسب سرعة \omterm{def:b1:nucleophilic-substitution:sn1}{التحلل بالمذيب}
SN1.
\end{exercise}

\begin{exercise}[$\star\star$]\label{exo:b1:nucleophilic-substitution:9}
يخضع بروميد ثانوي نقي ضوئيًا \omterm{def:b1:nucleophilic-substitution:sn1}{لتحلل بالمذيب}؛ فيدوّر الناتج
الضوء بنسبة \qty{12}{\percent} من القيمة المتوقعة لانقلاب
تام. أعطِ نسبتي \omterm{def:b1:stereochemistry-in-depth:enantiomers}{المتماكبين المرآويين} في الناتج
وفسّر.
\end{exercise}

\begin{exercise}[$\star\star\star$]\label{exo:b1:nucleophilic-substitution:10}
استنتج \omterm{def:b1:rate-laws:order}{قانون السرعة} العام \omterm{def:b1:nucleophilic-substitution:sn1}{للآلية SN1} \omterm{def:b1:elementary-steps:qssa}{بتقريب الحالة المستقرة}،
وبيّن أن إضافة أيون \omterm{def:b1:electronic-effects:nucleophile}{المجموعة المغادرة} \ce{Y-} في البداية تبطئ
التفاعل (\omterm{def:b1:precipitation:common-ion}{تأثير الأيون المشترك} في الحركية). وبأي شرط يكون
هذا التأثير مهملًا؟
\end{exercise}

\begin{exercise}[$\star\star\star$]\label{exo:b1:nucleophilic-substitution:11}
زمنا نصف تفاعل للبروموميثان مع الهيدروكسيد: مع
$[\ce{OH-}]_0 = [\ce{CH3Br}]_0 = C_0$، بيّن أن $1/[\ce{CH3Br}] = 1/C_0 +
kt$ وأعطِ $t_{1/2}$. ومع $[\ce{OH-}]_0 = 50\,C_0$، بيّن أن
التناقص أسي وأعطِ $t_{1/2}$. معطيات التمرين: $k =
\qty{2.0e-4}{L/(mol.s)}$، $C_0 = \qty{0.010}{mol/L}$.
\end{exercise}

\begin{exercise}[$\star\star\star$]\label{exo:b1:nucleophilic-substitution:12}
فسّر لماذا لا يتفاعل الكحول \ce{ROH} مع بروميد الصوديوم، لكنه
يتفاعل مع حمض الهيدروبروميك المركّز ليعطي \ce{RBr}. واكتب
الآلية لكحول ثالثي ولكحول أولي.
\end{exercise}

\section{مسألة: التحلل بالمذيب لكلوريد ثالثي البوتيل}

\begin{problem}[{مسألة نهاية الأسبوع --- الرتبة وثابت السرعة من
معطيات الموصلية، وزمن نصف التفاعل، والآلية وكيميائها الفراغية،
وطاقة تنشيط التحلل بالمذيب}]
\label{pb:b1:nucleophilic-substitution:1}
يُذاب 2-كلورو-2-ميثيل البروبان (كلوريد ثالثي البوتيل)، بتركيز
صغير، في خليط من الماء والبروبانون. وينتج تحلله بالمذيب
\ce{(CH3)3CCl + H2O -> (CH3)3COH + H+ + Cl-} أيونات، وتُسجَّل
موصلية المحلول $\kappa$، المتناسبة مع
كمية الأيونات المتكوّنة. وبعد زمن طويل جدًا تبلغ
$\kappa_\infty = \qty{2.000}{mS/cm}$ عند درجتي الحرارة كلتيهما. القياسات
(\unit{mS/cm}):

\begin{center}
\begin{tabular}{lcccccccc}
\hline
$t$ (min) & 0 & 5 & 10 & 15 & 20 & 30 & 45 & 60 \\
\hline
$\kappa$ عند \qty{25.0}{\celsius} & 0 & 0.172 & 0.329 & 0.473 & 0.605 & 0.835 & 1.110 & 1.321 \\
$\kappa$ عند \qty{35.0}{\celsius} & 0 & 0.507 & 0.885 & 1.168 & 1.379 & 1.654 & 1.856 & 1.940 \\
\hline
\end{tabular}
\end{center}
% data generated in figdata/bachelor-1/solvolysis.py (story data)

\medskip
\noindent\textbf{الجزء الأول --- الرتبة.}
\begin{enumerate}
  \item لماذا كانت الموصلية متناسبة مع كمية \omterm{def:b1:nucleophilic-substitution:substitution}{الركيزة}
        التي تفاعلت؟
  \item عبّر عن $[\mathrm{RCl}]/[\mathrm{RCl}]_0$ بدلالة $\kappa$
        و$\kappa_\infty$.
  \item اكتب القانون المتكامل لتفاعل من الرتبة الأولى
        والمقدار الذي يُرسم بدلالة $t$ لاختباره.
  \item احسب $\ln(\kappa_\infty - \kappa)$ عند \qty{25.0}{\celsius} لكل
        لحظة.
  \item هل المنحنى خط مستقيم؟ استنتج الرتبة.
  \item الماء بفائض كبير. فماذا يُخفي ذلك في الرتبة؟
\end{enumerate}

\medskip
\noindent\textbf{الجزء الثاني --- \omterm{def:b1:rate-laws:order}{ثابت السرعة} \omterm{def:b1:rate-laws:half-life}{وزمن نصف التفاعل}.}
\begin{enumerate}[resume]
  \item أوجد \omterm{def:b1:rate-laws:order}{ثابت السرعة} عند \qty{25.0}{\celsius}، بالوحدة \unit{s^{-1}}.
  \item احسب \omterm{def:b1:rate-laws:half-life}{زمن نصف التفاعل}.
  \item في أي لحظة يُستهلك \qty{90}{\percent} من \omterm{def:b1:nucleophilic-substitution:substitution}{الركيزة}؟
  \item أوجد \omterm{def:b1:rate-laws:order}{ثابت السرعة} عند \qty{35.0}{\celsius}.
  \item بأي عامل ازدادت السرعة لعشر درجات؟
  \item تحقق من قيمة واحدة من سلسلة \qty{35}{\celsius} بمقارنتها بالقانون.
\end{enumerate}

\medskip
\noindent\textbf{الجزء الثالث --- الآلية.}
\begin{enumerate}[resume]
  \item أي آلية توحي بها \omterm{def:b1:nucleophilic-substitution:substitution}{الركيزة} والمذيب؟
  \item اكتبها \omterm{def:b1:electronic-effects:curly-arrow}{بالأسهم المنحنية}.
  \item بيّن أنها توافق الرتبة التي وُجدت.
  \item إضافة هيدروكسيد الصوديوم بتركيز \qty{0.01}{mol/L} لا تكاد تغيّر
        \omterm{def:b1:rate-laws:rate}{سرعة اختفاء} الكلوريد. لماذا؟
  \item التجربة نفسها مع الكلوريد الثالثي \omterm{def:b1:stereochemistry-in-depth:stereocentre}{الكيرالي}
        3-كلورو-3-ميثيل الهكسان، النقي ضوئيًا، تعطي كحولًا نشاطه الضوئي
        يكاد يكون منعدمًا. فسّر.
  \item ارسم منحنى طاقة للتفاعل وعلّم
        \omterm{def:b1:elementary-steps:rds}{الخطوة المحددة للسرعة}.
  \item لماذا لا يكاد 1-كلوروبوتان يتفاعل في الظروف نفسها؟
\end{enumerate}

\medskip
\noindent\textbf{الجزء الرابع --- \omterm{def:b1:rate-laws:activation-energy}{طاقة التنشيط}.}
\begin{enumerate}[resume]
  \item اكتب \omterm{thm:b1:rate-laws:arrhenius}{قانون أرهينيوس}.
  \item عبّر عن $E_a$ بدلالة \omterm{def:b1:rate-laws:order}{ثابتي السرعة}.
  \item احسب $E_a$ ($R = \qty{8.314}{J/(K.mol)}$).
  \item ماذا تقيس هذه الطاقة في الآلية؟
  \item توقّع \omterm{def:b1:rate-laws:order}{ثابت السرعة} عند \qty{45.0}{\celsius}.
  \item اذكر \omterm{def:b1:rate-laws:activation-energy}{طاقة تنشيط} \omterm{def:b1:nucleophilic-substitution:sn1}{التحلل بالمذيب}، برقمين
        معنويين.
\end{enumerate}
\end{problem}
